\subsection{Collision-affinity bound}
\label{app:lower-proof-affinity}

We combine the midpoint identity with the one-edge Hellinger loss to
convert unresolved collisions into a multiplicative affinity penalty.
Recall our normalization: for probability masses \(p,q\) on the same
finite set \(\mathcal X\),
\(\HellingerAffinity(p,q):=\sum_{x\in\mathcal X}\sqrt{p(x)q(x)}\).
The squared Hellinger distance is
\(\sqHellinger(p,q):=\tfrac12\sum_{x\in\mathcal X}
(\sqrt{p(x)}-\sqrt{q(x)})^2=1-\HellingerAffinity(p,q)\).

\paragraph{Intuition.}
Each unresolved edge whose endpoints are committed together contributes
the factor $1-\HardEdgeDistance$.  Since
$\HardEdgeDistance\simeq\EdgeCoupling^2$,
\[
 1-(1-\HardEdgeDistance)^{\TotalCollisions}
 \simeq
 \min\{1,\EdgeCoupling^2\TotalCollisions\}.
\]
For justification, $1-(1-d)^z\le\min\{1,zd\}$ for integer $z\ge0$, while
$1-(1-d)^z\ge1-e^{-zd}\ge(1-e^{-1})\min\{1,zd\}$.
This is an edge-penalty scale; the next lemma places its expectation below
the actual squared-Hellinger risk.

\begin{lemma}[Affinity loss from unresolved collisions]
\label{app:lem:collision-affinity}
For every fixed decision-rule seed \(w\),
\begin{equation}
 \mathbb E_{I\sim\HardInstancePrior}
 \sqHellinger(\HardTargetLaw{I},\HardOutputLaw{I}{w})
 \ge
 \mathbb E_{I,H}
 \left[1-(1-\HardEdgeDistance)^{\TotalCollisions}\right],
 \label{app:eq:global-collision-affinity}
\end{equation}
where the expectation on the right is under
\eqref{app:eq:joint-instance-midpoint-law}.
\end{lemma}

\begin{proof}
\emph{1. Specify the edge components at a fixed history.}
Fix \(I,w\) and a reachable pre-commit history \(h\), and write
\(B_t=B_t(h)\). As in Step~4 of \cref{app:lem:unresolved-invariant}, the
two component laws on \(\Vocab^{e\cap B_t}\), for
\(e\in\HiddenMatching\), are
\begin{align*}
 p_e(z_{e\cap B_t})
 &:=\HardEdgeLaw{e}{I}
        (z_{e\cap B_t}\mid x_{e\cap G(h)}),\\
 \widehat p_e(z_{e\cap B_t})
 &:=\prod_{j\in e\cap B_t}
       \HardOracleRowMass{I}{y_t(h)}{j}{z_j}.
\end{align*}
The first line is the conditional probability of the current edge
coordinates given the already committed coordinates of that edge.
Both masses sum to one; on an empty coordinate set both equal one on the
empty assignment. The conditional product calculation in that step gives
\[
 \HardExactBatchLaw{t}{I}(z\mid h)
   =\prod_{e\in\HiddenMatching}p_e(z_{e\cap B_t}),
 \qquad
 \HardProductBatchLaw{t}{I}(z\mid h)
   =\prod_{e\in\HiddenMatching}\widehat p_e(z_{e\cap B_t}).
\]
These products are at fixed \(I,h\); each component is a block of batch
coordinates from one edge, not necessarily a single vertex.

\emph{2. Compute the factor of an unresolved collision.}
Partition the unresolved edges according to the number of selected endpoints:
\begin{equation}
 \mathcal E_t^{(r)}:=
 \{e\in\OpUnresolvedMatching{s_t}:
       |e\cap\ActiveBatch{t}|=r\},
 \qquad r\in\{0,1,2\}.
 \label{app:eq:commit-edge-partition}
\end{equation}
Thus the collision-edge set and its size are
\[
 \mathcal E_t^{(2)}
 =\{e\in\OpUnresolvedMatching{s_t}:
                       e\subseteq\ActiveBatch{t}\},
 \qquad
 |\mathcal E_t^{(2)}|=\CollisionCount{t},
 \qquad
 \ActiveBatch{t}=B_t\cap\UnresolvedVertices{t}.
\]
For \(e=\{i,j\}\in\mathcal E_t^{(2)}\), both endpoints belong to \(B_t\)
and neither belongs to \(G(h)\). Thus the exact component has no
conditioning within \(e\), while both oracle rows see their mate masked:
\[
 p_e(z_i,z_j)=\HardEdgeLaw{e}{I}(z_i,z_j),
 \qquad
 \widehat p_e(z_i,z_j)
 =\HardMarginal(z_i)\HardMarginal(z_j)
 =\HardIndependentEdgeLaw{e}(z_i,z_j).
\]
Here the oracle equality uses \eqref{app:eq:hard-oracle}, and the absence
of conditioning from other edges uses the target's edgewise product law.
Consequently,
\(\HellingerAffinity(p_e,\widehat p_e)
 =1-\sqHellinger(\HardEdgeLaw{e}{I},\HardIndependentEdgeLaw{e})
 =1-\HardEdgeDistance\), where \(\HardEdgeDistance\) is defined in
\eqref{app:eq:hard-edge-distance}. In terms of the actual edge masses,
this is
\begin{equation}
 \sum_{z_i,z_j}
 \sqrt{\HardEdgeLaw{e}{I}(z_i,z_j)
       \HardIndependentEdgeLaw{e}(z_i,z_j)}
 =1-\HardEdgeDistance.
 \label{app:eq:one-collision-affinity}
\end{equation}
\emph{3. Bound every remaining component and apply product affinity.}
Since \(p_e\) and \(\widehat p_e\) are probability masses,
\(0\le\HellingerAffinity(p_e,\widehat p_e)\le1\) by Cauchy--Schwarz.
This covers all noncollision components, including previously retired
edges where the oracle need not be exact. More specifically, an unresolved
crossing edge has \(p_e=\widehat p_e=\HardMarginal\), and an edge with
\(e\cap B_t=\varnothing\) has both masses equal to one on the empty
assignment; these components have affinity exactly one.

Apply \cref{app:lem:finite-hellinger-calculus} with index set
\(J=\HiddenMatching\), spaces \(\mathcal X_e=\Vocab^{e\cap B_t}\),
and component laws \(p_e,\widehat p_e\).
The two product laws to which it applies are exactly
\(\HardExactBatchLaw{t}{I}(\cdot\mid h)\) and
\(\HardProductBatchLaw{t}{I}(\cdot\mid h)\) from Step~1.
By the definition \eqref{app:eq:hard-local-affinity} and that lemma,
\begin{align*}
 \HardLocalAffinity{t}{I,h}
 &=\HellingerAffinity\!\left(
      \HardExactBatchLaw{t}{I}(\cdot\mid h),
      \HardProductBatchLaw{t}{I}(\cdot\mid h)\right)\\
 &=\prod_{e\in\HiddenMatching}\HellingerAffinity(p_e,\widehat p_e)\\
 &=\left(\prod_{e\in\mathcal E_t^{(2)}}(1-\HardEdgeDistance)\right)
   \left(\prod_{e\in\HiddenMatching\setminus\mathcal E_t^{(2)}}
                \HellingerAffinity(p_e,\widehat p_e)\right)\\
 &=(1-\HardEdgeDistance)^{\CollisionCount{t}}
   \prod_{e\in\HiddenMatching\setminus\mathcal E_t^{(2)}}
                \HellingerAffinity(p_e,\widehat p_e).
\end{align*}
Every factor in the remaining product lies in \([0,1]\), so
\begin{equation}
 \HardLocalAffinity{t}{I,h}
 \le(1-\HardEdgeDistance)^{\CollisionCount{t}}.
 \label{app:eq:local-collision-affinity}
\end{equation}
\emph{4. Multiply along the adaptive path and average.}
The preceding bound holds at each reached history, including histories
under the midpoint law. Multiplying over rounds and using
\(\TotalCollisions=\sum_t\CollisionCount{t}\) yields
\[
 \prod_t\HardLocalAffinity{t}{I,H_{t-1}}
 \le(1-\HardEdgeDistance)^{\TotalCollisions}.
\]
Using the adaptive identity \eqref{app:eq:adaptive-midpoint-identity}
first at fixed \(I,w\) and then averaging over \(I\) gives
\begin{align*}
 &\mathbb E_I\sqHellinger(\HardTargetLaw{I},\HardOutputLaw{I}{w})\\
 &\quad=1-\mathbb E_I
    \mathbb E_{H\sim\HardMidpointLaw{I}{w}}
       \prod_t\HardLocalAffinity{t}{I,H_{t-1}}\\
 &\quad\ge1-\mathbb E_{I,H}(1-\HardEdgeDistance)^{\TotalCollisions}\\
 &\quad=\mathbb E_{I,H}
       [1-(1-\HardEdgeDistance)^{\TotalCollisions}],
\end{align*}
where the joint expectation uses the conditional midpoint law given \(I,w\).
This proves \eqref{app:eq:global-collision-affinity}.
\end{proof}
